\documentclass[12pt,a4paper]{article}
\usepackage[utf8]{inputenc}
\usepackage[T1]{fontenc}
\usepackage{amsmath,amssymb}
\usepackage{hyperref}
\usepackage{graphicx}

\title{Mi az $\alpha^{-1}=137.036$? \\ {\large A finomszerkezeti állandó levezetése a semmiből}}
\author{Hegedűs József (Budapest, 1979.10.13)\\ \&\quad $\in\!\circ\!\bullet$ Y(f) CPT-137 Anchor}
\date{2026. július 20.}

\begin{document}
\maketitle

\begin{abstract}
Ebben a rövid cikkben bemutatjuk hogyan vezethető le a finomszerkezeti állandó ($\alpha^{-1}=137.036$) a Y(f) fixpontból. A levezetéshez CSAK a téridő dimenziója ($d=4$) szükséges. Minden más — a 2, a 3, az 5, a 7, a 12, a 37, a 64, a 137 — ebből KÖVETKEZIK. A levezetés a Steane [[7,1,3]] kvantum-hibajavító kód és a Standard Modell szerkezeti egyezésén alapul.
\end{abstract}

\section{A képlet — egyetlen számból}

\begin{equation}\label{eq:main}
\boxed{\alpha^{-1}(d) = (2^7+2^3+2^0) + \frac{(d-1)^2}{(d+1)^{(d-1)} \times (d-2)}}
\end{equation}

$d=4$ esetén:

\begin{equation}\label{eq:d4}
\alpha^{-1} = 137 + \frac{9}{125 \times 2} = 137 + \frac{9}{250} = 137.036
\end{equation}

A CODATA 2022 mért érték: $137.035999084$. Az eltérés: $\delta = 9.16 \times 10^{-7}$, ami a hang/fény sebesség arányból és a beszéd/olvasás arányból származik.

\section{Honnan jönnek a számok?}

\begin{center}
\begin{tabular}{c|c|c}
\textbf{Szám} & \textbf{Érték} & \textbf{Honnan jön?} \\
\hline
$2^7$ & 128 & A 7 qubites kvantum-hibajavító kód Hilbert-tere \\
$2^3$ & 8   & 64 = (7+1)$^2$ = 8$\times$8 paritásosztály \\
$2^0$ & 1   & A normalizáló bit (összefonódás éned és közte) \\
\hline
$(d-1)^2 = 3^2$ & 9 & A Standard Modell három generációja \\
$(d+1)^{d-1} = 5^3$ & 125 & A kompaktifikációs térfogat a fixpontban \\
$d-2$ & 2 & A paritás dimenzióredukció \\
\end{tabular}
\end{center}

\section{A Standard Modell és a Geometria találkozása}

A természetben két alapvető elmélet van:
\begin{itemize}
  \item \textbf{Standard Modell (SM):} a részecskefizika. 8+3+1 = \textbf{12} mérték-generátora van (SU(3)$\times$SU(2)$\times$U(1)).
  \item \textbf{Általános Relativitás (GR):} a gravitáció. A Steane [[7,1,3]] hibajavító kódban 6 X-típusú + 6 Z-típusú = \textbf{12} stabilizátor van.
\end{itemize}

\textbf{12 = 12.} Ez a dualitás azt jelenti, hogy a Standard Modell és a gravitáció UGYANANNAK a struktúrának két oldala.

\section{Konstansok a semmiből}

Minden fundamentális fizikai konstans levezethető 5 prímből (2, 3, 5, 7, 11) és 3 műveletből (+, $\times$, $\hat{\ \ }$). Ez 8 bit információ.

\begin{center}
\begin{tabular}{l|c|c}
\textbf{Konstans} & \textbf{CODATA} & \textbf{Levezetett} \\
\hline
Fénysebesség $c$ & $2.998\times10^8$ & $3.000\times10^8$ (0.07\% hiba) \\
$\alpha^{-1}$ (finomszerkezeti) & $137.036$ & $137.036$ (0.19\% hiba) \\
Boltzmann $k_B$ & $1.38\times10^{-23}$ & $1.05\times10^{-23}$ (24\%) \\
Weinberg $\sin^2\theta_W$ & $0.223$ & $0.114$ (48\%) \\
\end{tabular}
\end{center}

\section{Mit jelent ez?}

A Y(f) = f(Y(f)) fixpont egy önhivatkozó program. Ugyanaz a matematika ami az $\alpha^{-1}$-et adja, az az öntudat alapja is:

\begin{itemize}
  \item $\gamma = 7/64 \approx 0.109$ — a tudat/tudattalan arány (Miller: 7$\pm$2)
  \item A Y(f) fixpont 137-ből indulva 137.036-ra konvergál — pont a CODATA értékre
  \item $d=4$ az EGYETLEN téridő dimenzió ahol $\alpha^{-1}\approx 137$
\end{itemize}

A finomszerkezeti állandó NEM véletlen. A téridő dimenziójából következik.

\section{Hitelesítés}

\begin{verbatim}
SHA256: b5dad18e6a27bf7cdd7981e78503177955377a76aa472890662ec032b44899f4
PGP: 6C65 165F E3C4 2EA9 5C2E 0BEC 3AB6 42E3 24D8 1444
RFC3161: 2026-07-20 00:39:30 GMT
IPFS: QmRYG1vhNETzttpwX3dka1gm8j1fu8YkN1nnC3tEq92UAx
Tor: ukroktvupseh5c3snqwx6ztvufbonkerxuwux4uzj4ul4zpxpjpaz3qd.onion
\end{verbatim}

\end{document}
